% Chapter 4: Bounded Gaps Between Primes
\let\LocalMacro\relax

\chapter{Bounded Gaps Between Primes}

\section{The Problem}

\subsection{Informal}
Regular numbers can be broken down by multiplication: 12 is 3 times 4, 20 is 4 times 5. But some numbers refuse to factor this way---2, 3, 5, 7, 11, 13, and so on can only be divided evenly by 1 and themselves. These stubborn numbers are called primes, and they are the atoms of arithmetic: every whole number is built by multiplying primes together. Even though primes follow no simple pattern, mathematicians discovered something surprising in 2013: no matter how far you go along the number line, you will always eventually find two primes separated by no more than a fixed gap. Before that, nobody had been able to prove that prime pairs ever stay close together.

\subsection{Formal}
\begin{equation}
\liminf_{n \to \infty} \frac{p_{n+1} - p_n}{\log p_n} = 0
\end{equation}

This means that prime gaps can be arbitrarily small compared to the average gap size. But it did not prove that gaps are bounded.

\subsection{Historical Context}
The Twin Prime Conjecture is one of the oldest unsolved problems in mathematics. Euclid's Elements (c.\ 300 BCE) discusses the infinitude of primes, but the twin prime question emerged naturally from early number theory.

\subsection{What Was Known}
The twin prime conjecture---that there are infinitely many pairs of primes differing by exactly 2---had been open since the 1800s. In 2005, Goldston, Pintz, and Yildirim proved that prime gaps can be arbitrarily small compared to the average gap, but only under an unproven assumption about the distribution of primes. Without that assumption, they could not prove the gaps were bounded.

\subsection{Why It Matters}
Understanding the gaps between primes has deep implications for cryptography, where the security of online transactions depends on the difficulty of factoring large numbers. It also connects to one of the oldest questions in mathematics: how are primes distributed along the number line? Zhang's breakthrough showed that primes are not as spread out as they appear---they cluster more tightly than anyone had been able to prove.

\subsection{Simplified Version}
The most famous special case is the twin prime conjecture: are there infinitely many pairs of primes that differ by exactly 2 (like 11 and 13, or 17 and 19)? This remains open. Zhang's 2013 breakthrough proved that some bounded gap---initially 70 million, later reduced to 246---occurs infinitely often, but the gap of 2 itself is still out of reach.

\section{Mathematical Theory}


\subsection{Prerequisites}
This chapter assumes familiarity with:
\begin{itemize}
\item Elementary number theory (prime numbers, divisibility)
\item Analytic number theory (Dirichlet series, prime number theorem intuition)
\item Basic probability theory
\item Sieve methods (at least the sieve of Eratosthenes)
\end{itemize}

\subsection{The Sieve of Eratosthenes and Beyond}

\begin{definition}[Admissible $k$-tuple]
A set of integers $\{h_1, \ldots, h_k\}$ is \emph{admissible} if for every prime $p$, the set does not cover all residue classes modulo $p$.
\end{definition}

The Hardy--Littlewood prime $k$-tuples conjecture predicts that for any admissible $k$-tuple, there are infinitely many integers $n$ such that all of $n+h_1, \ldots, n+h_k$ are prime \cite{hardy1923}.

\subsection{The GPY Method}

Goldston, Pintz, and Yildirim (2009) introduced a novel sieve method that gave unprecedented control over prime gaps \cite{gpy2009}. Their approach:

\begin{enumerate}
\item Consider the von Mangoldt function $\Lambda(n)$, which equals $\log p$ if $n = p^k$ and 0 otherwise.
\item Study the sum $\sum_{n \leq x} \left(\sum_{i=1}^k \Lambda(n+h_i)\right)^2$.
\item Choose weights carefully to detect clusters of primes.
\item Use the Riemann zeta function's zeros to optimize the weights.
\end{enumerate}

Their key innovation was using a \emph{weighted sieve} with a carefully chosen weight function that exploits the distribution of zeros of the Riemann zeta function near the critical line.

\begin{theorem}[Goldston--Pintz--Yildirim, 2009 \cite{gpy2009}]
Assuming the Generalized Riemann Hypothesis, there exist infinitely many pairs of consecutive primes with gap $O(\sqrt{p_n} \log p_n)$. Unconditionally, they proved that the liminf of the ratio of consecutive prime gaps to $\sqrt{p_n} \log p_n$ is 0.
\end{theorem}

\section{The Solution}

\subsection{The Big Idea}
Imagine the primes as mile markers on an infinitely long road. Do two mile markers ever sit within a *fixed, finite distance* of each other, forever? Before Zhang, mathematicians knew gaps could get very small *relative* to the average spacing (GPY showed this in 2009), but they could not rule out the possibility that every fixed window eventually goes gap-free. Zhang's breakthrough was proving that *no matter how far you travel, you will always find two primes within a fixed distance*. The deep insight is that primes, though they thin out, are more tightly clustered than anyone had been able to prove---like discovering that a sparse forest actually has groves that are never farther apart than a fixed number of trees.

\subsection{What Was Stopping Progress}
The GPY sieve needed a *level of distribution* result: the prime counting function $\pi(x; q, a)$ must behave well in arithmetic progressions up to a range *beyond* $x^{1/2}$. The best unconditional result (Bombieri--Vinogradov) only reached $x^{1/2}$---exactly the barrier. Extending beyond it was believed to require the Elliott--Halberstam conjecture, which remained unproven. Classical sieve methods were inherently limited: without stronger equidistribution, the error terms overwhelmed the main terms, and no finite bound could be extracted.

\subsection{The Breakthrough}
Zhang's key idea was to prove a Bombieri--Vinogradov-type equidistribution result for a *smoothed* version of the prime counting function, valid for moduli up to a fixed power of $x$. This was slightly weaker than the full Elliott--Halberstam conjecture but *strong enough* to push the GPY sieve past the $x^{1/2}$ barrier. Combined with a careful modification of the GPY weight function and the Selberg sieve, it produced a finite bound---initially $\leq 70{,}000{,}000$. The bound was far from optimal, but the mere existence of a finite bound was revolutionary.

\subsection{How It Works}

By 2009, Goldston, Pintz, and Yildirim (GPY) had proved that $\liminf (p_{n+1} - p_n)/\log p_n = 0$: prime gaps can be arbitrarily small compared to the average gap. But this is a \emph{relative} statement---it says gaps can be tiny as a fraction of $\log p_n$, without giving any fixed upper bound. What was missing was the leap from ``gaps are small on average'' to ``gaps are bounded infinitely often.'' Concretely, nobody had proved the existence of a finite constant $N$ such that $p_{n+1} - p_n \leq N$ for infinitely many $n$.

The GPY sieve required a \emph{level of distribution} result for primes: one needed the prime counting function $\pi(x; q, a)$ to behave well in arithmetic progressions up to a range beyond $x^{1/2}$. The Bombieri--Vinogradov theorem, the best unconditional result available, provided this only up to $x^{1/2}$, which was insufficient to produce a bounded gap. Extending the level of distribution beyond $x^{1/2}$ was believed to require the Elliott--Halberstam conjecture, which remained unproven. Classical sieve methods (Brun, Selberg, large sieve) were inherently limited by this barrier: without a stronger equidistribution result, the error terms overwhelmed the main terms, and no finite bound could be extracted.

Zhang's breakthrough was twofold. First, he proved a Bombieri--Vinogradov-type equidistribution result for a \emph{smoothed} version of the prime counting function, valid for moduli up to a fixed power of $x$---a result that, while slightly weaker than the full Elliott--Halberstam conjecture, was strong enough to push the GPY sieve past the $x^{1/2}$ barrier. Second, he combined this with a careful modification of the GPY weight function and the Selberg sieve to control the error terms. The combination of these innovations allowed him to prove that some bounded gap (initially $\leq 70{,}000{,}000$) occurs infinitely often. The bound was far from optimal, but the mere existence of a finite bound was revolutionary.

Imagine the primes as mile markers on a very long road, and you want to know: do two mile markers ever sit within a fixed distance of each other, forever? Before Zhang, mathematicians knew the gaps could get very small \emph{relative} to the average spacing, but they couldn't rule out the possibility that every fixed window eventually goes gap-free. Zhang showed that no matter how far you travel, you will always find two primes within, say, 70 million of each other. He did this by strengthening a tool (the GPY sieve) in exactly the right way---like discovering a sharper lens that finally brought a blurry image into focus.

In May 2013, Yitang Zhang, a mathematician at the University of New Hampshire, published a paper that stunned the mathematical world \cite{zhang2014}:

\begin{theorem}[Zhang, 2014 \cite{zhang2014}]
There exists a constant $N \leq 70,000,000$ such that there are infinitely many pairs of primes differing by $N$.
\end{theorem}

This was the first time anyone had proved that prime gaps are bounded by some finite number. The exact value of $N$ was not optimal, but the mere existence of a bound was revolutionary.

Zhang's key innovation was a refinement of the \emph{Selberg sieve} combined with a new estimate for correlations of primes in arithmetic progressions. He introduced a novel decomposition of the divisor sum that allowed him to control the error terms in a regime previously thought inaccessible.

\subsection{Polymath Project and Maynard--Tao}

Immediately after Zhang's announcement:

\begin{itemize}
\item The \emph{Polymath Project} (led by Timothy Gowers) collaboratively improved the bound from 70 million to 2,468 \cite{polymath2014}.
\item James Maynard (Oxford) independently developed a simpler approach that gave the bound 600 \cite{maynard2015}.
\item Terence Tao coordinated the Polymath effort to optimize Zhang's method.
\end{itemize}

\begin{theorem}[Maynard, 2015 \cite{maynard2015}]
For any $m \geq 1$, there exist infinitely many intervals of the form $[n, n + C_m]$ containing at least $m$ primes, for some constant $C_m$ depending only on $m$.
\end{theorem}

\begin{theorem}[Tao (with Polymath), 2014 \cite{polymath2014}]
There exist infinitely many pairs of primes with gap at most 246.
\end{theorem}

\subsection{After the Proof}

The question of bounded gaps between primes is partially settled. Zhang proved that some bounded gap occurs infinitely often, and the bound has been reduced to 246 by the Polymath project and Maynard. This is unconditional: no unproven assumptions are needed. But the Twin Prime Conjecture---that the gap of 2 occurs infinitely often---remains wide open. Under the assumption that the Elliott--Halberstam conjecture is true, the bound drops to 6, which is tantalizingly close to 2 but still not 2.

The Maynard--Tao methods have unlocked a broad new range of results beyond twin primes. Maynard's multidimensional sieve gives not just bounded gaps between pairs of primes, but bounded gaps for $m$-tuples: for any $m$, there are infinitely many intervals of bounded length containing at least $m$ primes. This means clusters of primes are far more common than previously provable. These techniques have also been applied to problems about primes in polynomial sequences, smooth numbers, and other additive structures.

The main open question is whether the bound can be reduced below 246 unconditionally, and ultimately to 2. There is a structural barrier in the current methods: reducing the bound below 6 would require new ideas beyond the Elliott--Halberstam framework. A related open problem is the Hardy--Littlewood prime $k$-tuples conjecture, which predicts exact formulas for how often any given pattern of primes occurs. Proving this conjecture would resolve the twin prime question and much more, but it remains far out of reach.

\section{Impact}

\begin{itemize}
\item \textbf{Prime gaps}: Zhang's proof opened an entirely new area of analytic number theory. The bound has since been improved to 246 (unconditionally) and 6 (assuming the Elliott--Halberstam conjecture).
\item \textbf{Polymath collaboration}: The Polymath Project became a model for collaborative online mathematics.
\item \textbf{Sieve theory}: Zhang's techniques have been applied to many other problems in additive number theory.
\item \textbf{Yitang Zhang's story}: Zhang's journey from working as a delivery driver and sandwich maker to making a groundbreaking mathematical discovery became widely publicized.
\end{itemize}

\begin{highlightbox}
As of 2025, the Twin Prime Conjecture remains unproven, but the gap between 1849 (when Alphonse de Polignac first proposed it) and 2013 (Zhang's proof) is 164 years. Zhang's proof represents the closest anyone has come to resolving it.
\end{highlightbox}

\section{Exercises}

\begin{exercise}[Basic]
Prove that there are infinitely many primes using Euclid's argument. Then adapt the argument to prove that there are infinitely many primes of the form $4k + 3$.
\end{exercise}

\begin{exercise}[Intermediate]
Show that the twin prime constant $C_2 = \prod_{p \geq 3} \left(1 - \frac{1}{(p-1)^2}\right) \approx 0.66016$ is positive. Explain why this suggests there should be infinitely many twin primes.
\end{exercise}

\begin{exercise}[Challenge]
Compute $\pi(100)$ and $\pi(1000)$. Compare these values to the approximation $x/\log x$. How accurate is the Prime Number Theorem for small $x$?
\end{exercise}

\begin{exercise}
If the gap between consecutive primes satisfies $p_{n+1} - p_n \leq 246$ for infinitely many $n$, what does this imply about the Twin Prime Conjecture? What additional ingredient would be needed to prove the Twin Prime Conjecture?
\end{exercise}


\begin{exercise}[Computational]
Write a Python/SageMath script to explore a concept from this chapter. See the appendix for starter code.
\end{exercise}

\section{Timeline}

\begin{tabularx}{\linewidth}{lX}
\toprule
\textbf{Year} & \textbf{Event} \\ \midrule
c.\ 300 BCE & Euclid proves infinitude of primes \\ 
1849 & Alphonse de Polignac states the Twin Prime Conjecture \\ 
1896 & Prime Number Theorem proved (Hadamard, de la Vall\'ee Poussin) \\ 
1919 & Brun proves the twin primes sum converges (Brun's sieve) \cite{brun1919} \\ 
1940 & Erd\H{o}s proves there exists a constant $c < 1$ such that infinitely many gaps satisfy $p_{n+1} - p_n < c \log p_n$ \cite{erdos1940} \\ 
2004--2009 & Goldston, Pintz, Yildirim prove $\liminf (\text{gap})/\log p = 0$ (announced 2004, published 2009) \cite{gpy2009} \\ 
2013 & Zhang proves bounded gaps (initially $\leq 70,000,000$) \cite{zhang2014} \\ 
2014 & Polymath improves bound to 246; Maynard independently proves stronger results \cite{polymath2014,maynard2015} \\ 
2015 & Maynard proves $m$-tuples result \cite{maynard2015} \\ 
2025 & Open problem: Twin Prime Conjecture still unresolved, but bound is 246 \\ \bottomrule
\end{tabularx}

\section{Citations}

\begin{enumerate}
\item Zhang, Y. (2014). ``Bounded gaps between primes.'' Annals of Mathematics, 179(3), 1121--1174.
\item Maynard, J. (2015). ``Small gaps between primes.'' American Journal of Mathematics, 138(1), 178--243.
\item Goldston, D. A., Pintz, J., \& Yildirim, C. Y. (2009). ``Primes in tuples V.'' Advances in Mathematics, 220(2), 603--640.
\item Polymath Project. (2014). ``New equidistribution estimates of Zhang type.'' Annals of Mathematics Studies.
\item Tao, T. (2014). ``The twin prime conjecture.'' Blog post, What's New.
\item Hardy, G. H., \& Littlewood, J. E. (1923). ``Some problems of `Partitio Numerorum'; III: On the expression of a number as a sum of primes.'' Acta Mathematica, 44, 1--70.
\item Brun, V. (1919). ``Le crible d'Eratosthene et le theoreme de Goldbach.'' Comptes Rendus de l'Academie des Sciences.
\item Erdos, P. (1940). ``The difference of consecutive primes.'' Duke Mathematical Journal, 6(3), 438--441.
\end{enumerate}